\section{The split form and binding in neutral signature}
\label{sec:neutral}

The norm in \eqref{eq:alpha} is positive definite, and that is what makes $g^{\!\top}g\succeq0$ in \cref{prop:inverted} and $|p|^2\ge0$ in
\cref{thm:virial}. But the arena does not hand us a positive form; it hands us $\eta$, of signature $(3,3)$. Redoing the derivation with
$\eta$ changes one inequality, and with it the verdict.

\subsection{The split-form Lagrangian}

Write $G^\dagger:=\eta^{-1}G^{\!\top}\eta$ for the $\eta$-adjoint and take
\begin{equation}
\label{eq:splitL}
\mathcal L=\frac m2\,\eta\bigl(\alpha_+X,\alpha_+X\bigr),\qquad p=m\,\eta\,(\alpha_+X),\qquad \pi=\eta^{-1}p,\qquad \dot p=\kappa\,G^{\!\top}p ,
\end{equation}
with second-order form and energy
\begin{equation}
\label{eq:splitEg}
\ddot X+2\kappa a\dot X-\kappa^2G^\dagger GX=0,\qquad a=\tfrac12(G-G^\dagger),\qquad E=\frac m2\eta(\dot X,\dot X)-\frac{m\kappa^2}2\eta(GX,GX).
\end{equation}
Everything in \cref{sec:shifted} survives with the transpose replaced by the $\eta$-adjoint: factorisation holds exactly when $a=0$, and $E$
is conserved because $a$ is $\eta$-antisymmetric. What does not survive is the sign of the virial identity.

\begin{theorem}[The split virial identity]
\label{thm:splitvirial}
Let $X$ be a section along a path in a manifold with metric of signature $(3,3)$ and let $D/D\tau$ come from any metric-compatible connection,
torsion included. Then for every $G$,
\begin{equation}
\label{eq:splitvirial}
\frac{d}{d\tau}\,\eta(X,\pi)=\frac{\eta(\pi,\pi)}m .
\end{equation}
Consequently every closed orbit satisfies $\oint\eta(\pi,\pi)\,d\tau=0$: a bound orbit has null momentum on average.
\end{theorem}

\begin{proof}
Note that $\eta(X,\pi)=\langle p,X\rangle$. Metric compatibility gives $\tfrac{d}{d\tau}\eta(X,\pi)=\eta(DX/D\tau,\pi)+\eta(X,D\pi/D\tau)$, and torsion does
not enter this Leibniz rule. Substituting $DX/D\tau=\pi/m-\kappa GX$ and $D\pi/D\tau=\kappa G^\dagger\pi$ gives $\eta(\pi,\pi)/m-\kappa\,\eta(GX,\pi)+\kappa\,\eta(X,G^\dagger\pi)$,
and the last two terms cancel by the definition of the adjoint.
\end{proof}

Neither curvature nor torsion appears in \eqref{eq:splitvirial}; the identity sees only $\eta$. In definite signature $\eta(\pi,\pi)$ has
one sign and the identity forbids recurrence, which is \cref{thm:virial}. In split signature it does not, and the question becomes whether
solutions with null momentum on average exist. They do, and the operator that produces them is one we already have.

\subsection{Binding exactly in neutral signature}

\begin{theorem}[Factorisation and binding coexist exactly in neutral signature]
\label{thm:neutral}
Let $\eta$ have signature $(P,Q)$ on $\R^N$. There is a linear map $G$ with
\begin{enumerate}[label=(\roman*),leftmargin=2.2em]
\item $G=G^\dagger$, so that the factorisation of \cref{thm:factor} holds, and
\item $G^2=-\id$, so that \eqref{eq:splitEg} reads $\ddot X=-\kappa^2X$ and every solution is periodic of period $2\pi/\kappa$,
\end{enumerate}
if and only if $P=Q$. Such a $G$ is an anti-isometry, $\eta(GX,GY)=-\eta(X,Y)$, and on $\R^{3,3}$ one of them, maximal in that it uses every
paired direction, is the binding root of \cref{prop:twoK}:
\begin{equation}
\label{eq:Gpair}
G:\ e_i\longmapsto e_{i+3},\qquad e_{i+3}\longmapsto-e_i\qquad(i=1,2,3).
\end{equation}
\end{theorem}

\begin{proof}
Given (i) and (ii), $\eta(GX,GY)=\eta(X,G^\dagger GY)=\eta(X,G^2Y)=-\eta(X,Y)$, so $G$ carries a maximal positive subspace isometrically onto a
negative one and back, which forces $P=Q$. Conversely \eqref{eq:Gpair} satisfies $G^2=-\id$, and
$\eta(Ge_i,e_{i+3})=\eta(e_{i+3},e_{i+3})=-1=\eta(e_i,-e_i)=\eta(e_i,Ge_{i+3})$, so $G=G^\dagger$. With $G^\dagger G=G^2=-\id$, \eqref{eq:splitEg}
has $a=0$ and becomes $\ddot X=-\kappa^2X$.
\end{proof}

This reverses the verdict of \cref{sec:shifted} at its root. In definite signature the factorised branch is an inverted oscillator and
nothing binds; in neutral signature the same factorisation condition $G=G^\dagger$ admits operators of purely imaginary spectrum, and the
maximal one makes the whole interior a harmonic well of one frequency. The obstruction of \cref{thm:virial} was a theorem about the
positive form, not about the arena. Two qualifications keep the statement honest. The normalisation $G^2=-\id$ imposes one common
frequency; any $\eta$-self-adjoint shift built from independent paired blocks $\mu_iJ_2$ with distinct $\mu_i>0$ also factorises and binds, at
unequal frequencies, so the construction is a family and $G^2=-\id$ picks a member. And no equation of motion for $G$ has been written:
$G$ is adopted, and calling its consequences a selection would be premature.

\begin{proposition}[How much can bind]
\label{prop:howmuch}
For $G=G^\dagger$ an eigenvector $v$ with non-real eigenvalue satisfies $\eta(v,\bar v)=0$, so each conjugate pair $\pm i\mu$ spans an
$\eta$-neutral two-plane. The oscillating subspace has signature $(k,k)$ with $k\le\min(P,Q)$. On Minkowski space $\R^{1,3}$ at most a
single two-plane can bind; on $\R^{3,3}$ all six directions can, simultaneously and at one frequency.
\end{proposition}

\begin{example}[One pairing plane]
\label{ex:oneplane}
In the plane $(e_1,e_4)$ write $X=x\,e_1+y\,e_4$, so that $G$ acts as $\bigl(\begin{smallmatrix}0&-1\\1&0\end{smallmatrix}\bigr)$ and $z=x+iy$. The
general solution of $\ddot X=-\kappa^2X$ is $z(\tau)=c_-e^{-i\kappa\tau}+c_+e^{i\kappa\tau}$, an ellipse with semi-axes $|c_-|+|c_+|$ and
$\bigl||c_-|-|c_+|\bigr|$. The momentum $\alpha_+X=\dot X+\kappa GX$ corresponds to $\dot z+i\kappa z=2i\kappa\,c_+e^{i\kappa\tau}$, so it is carried
entirely by the counter-rotating part, and the zero-momentum orbits are the circles $z=c_-e^{-i\kappa\tau}$, traversed at angular velocity
$-\kappa$. The virial quantity $\eta(\pi,\pi)=m^2\eta(\alpha_+X,\alpha_+X)=m^2\Re\bigl((2i\kappa c_+e^{i\kappa\tau})^2\bigr)$ oscillates with mean
zero, as \cref{thm:splitvirial} requires of a closed orbit.
\end{example}

\begin{figure}[t]
\centering
\begin{tikzpicture}
\begin{axis}[name=left,width=6.8cm,height=6.8cm,axis equal image,xmin=-5.4,xmax=5.4,ymin=-5.4,ymax=5.4,
  axis lines=middle,xtick=\empty,ytick=\empty,xlabel={$x$},ylabel={$y$},
  title={\small positive form, $g=J$}]
\addplot[gray,dashed,thick,domain=0:360,samples=90] ({cos(x)},{sin(x)});
\addplot[scol,thick,domain=0:5*pi,samples=600,variable=\t] ({cos(deg(\t))+0.3*\t*sin(deg(\t))},{-sin(deg(\t))+0.3*\t*cos(deg(\t))});
\node[scol,anchor=north west,font=\small] at (axis cs:-5.3,5.3) {$p_0\neq0$: spiral};
\node[gray,anchor=south west,font=\small] at (axis cs:-5.3,-5.3) {dashed: $p=0$};
\end{axis}
\begin{axis}[at={(left.east)},anchor=west,xshift=1.0cm,width=6.8cm,height=6.8cm,axis equal image,xmin=-2.0,xmax=2.0,ymin=-2.0,ymax=2.0,
  axis lines=middle,xtick=\empty,ytick=\empty,xlabel={$x$},ylabel={$y$},
  title={\small split form, $G$ the pairing}]
\addplot[gray,dashed,thick,domain=0:360,samples=90] ({cos(x)},{sin(x)});
\addplot[tcol,thick,domain=0:360,samples=180] ({1.4*cos(x)},{-0.6*sin(x)});
\addplot[gcol,very thick,domain=0:360,samples=180] ({0.85*cos(x)-0.85*sin(x)},{0.85*cos(x)-0.85*sin(x)});
\node[tcol,anchor=north east,font=\small] at (axis cs:1.95,-0.7) {$|c_+|<|c_-|$};
\node[gcol,anchor=north west,font=\small] at (axis cs:-1.95,-1.45) {$|c_+|=|c_-|$};
\node[gray,anchor=south east,font=\small] at (axis cs:-0.55,0.95) {$p=0$};
\end{axis}
\end{tikzpicture}
\caption{Interior motion in one pairing plane (\cref{ex:oneplane}). Left: with the positive form and an antisymmetric shift the
zero-momentum orbits are circles, and any momentum turns them into outward spirals, as \cref{thm:virial} requires. Right: with the split form
and the binding root $G$ every orbit is closed; the zero-momentum orbits are circles, the momentum-carrying ones ellipses, degenerating to a
segment when the two rotating parts have equal size.}
\label{fig:interior}
\end{figure}

\subsection{The orbits as duality rotations}

Read through the field-strength dictionary of \cref{thm:emfield}, an interior configuration $X$ is a two-form and the binding root is the
Hodge star. The bound motion then has a transparent meaning.

\begin{proposition}[Bound orbits in field-strength language]
\label{prop:orbitsEM}
Let $G$ be the pairing \eqref{eq:Gpair} and $z=\bB+i\bE$.
\begin{enumerate}[label=(\roman*),leftmargin=2.2em]
\item The zero-momentum orbits are $X(\tau)=e^{-\kappa G\tau}X_0$, that is $z(\tau)=e^{-i\kappa\tau}z_0$: continuous duality rotations, along which
$q(z(\tau))=e^{-2i\kappa\tau}q(z_0)$. A non-null orbit is a single spacetime plane at exactly four instants per period, alternately of magnetic
and electric type, and a null orbit is a single plane at every instant.
\item Every solution is $z(\tau)=c_-e^{-i\kappa\tau}+c_+e^{i\kappa\tau}$ with $c_\pm\in\C^3$, the momentum corresponds to $2i\kappa c_+e^{i\kappa\tau}$, and the
period averages are
\[
\langle q\rangle=2\,c_-\cdot c_+,\qquad \langle\bE\cdot\bB\rangle=\Im(c_-\cdot c_+),\qquad \langle|\bB|^2-|\bE|^2\rangle=2\Re(c_-\cdot c_+).
\]
\item The energy is $\mathcal E=2m\kappa^2\Re(c_-\cdot c_+)$, and the complex quantity $\tfrac m2\bigl(q(\dot z)+\kappa^2q(z)\bigr)=2m\kappa^2\,c_-\cdot c_+$ is
conserved; its imaginary part is a second constant of motion.
\item On the orbits with $c_+=\bar c_-$ the vector $z$ is real at every instant: $\bE=0$ and $\bE\cdot\bB=0$ throughout, while the momentum is not
zero.
\end{enumerate}
\end{proposition}

On a zero-momentum orbit $\bE\cdot\bB(\tau)=\tfrac12|q_0|\sin(\phi_0-2\kappa\tau)$ with $q_0=|q_0|e^{i\phi_0}$, and the rest-frame field of
\cref{sec:orth} turns from purely magnetic through screws to purely electric and back, at the fixed rest-frame energy density
$\tfrac12|q_0|$. So the binding flow is what breaks orthogonality on every non-null orbit without momentum, while orbits that carry momentum
can stay orthogonal throughout. The mean non-orthogonality and the energy of a general orbit are the imaginary and real parts of one conserved
complex number $c_-\cdot c_+$, and neither determines the other.

\begin{theorem}[The parity-odd term is a total derivative]
\label{thm:EBsplit}
For the split-form Lagrangian with the pairing,
\[
\mathcal L_\kappa=\frac m2\Bigl(\eta(\dot X,\dot X)-\kappa^2\eta(X,X)\Bigr)-m\kappa\,\frac{d}{d\tau}\bigl(\bE\cdot\bB\bigr).
\]
Consequently the cross term has identically vanishing Euler--Lagrange expression, the equation $\ddot X=-\kappa^2X$ comes from the first bracket
alone, and the momenta of $\mathcal L_\kappa$ and $\mathcal L_{-\kappa}$ differ by $2m\kappa\,\eta GX$. Parity maps $\mathcal L_\kappa$ to $\mathcal L_{-\kappa}$
and exchanges the orbits without momentum with those carried entirely by the counter-rotating part, while $C$ and $T$ leave $\mathcal L_\kappa$
invariant.
\end{theorem}

\begin{proof}
$G$ is $\eta$-self-adjoint, so $\tfrac{d}{d\tau}\eta(X,GX)=2\eta(\dot X,GX)$, and $\eta(X,GX)=-2\,\bE\cdot\bB$ by \cref{thm:emfield}; $\eta(GX,GX)=-\eta(X,X)$
because $G$ is an anti-isometry. Expanding $\eta(\dot X+\kappa GX,\dot X+\kappa GX)$ gives the formula. Parity reverses $\bE$ and fixes $\bB$, so
$PGP^{-1}=-G$ and $P$ acts on $z$ as complex conjugation.
\end{proof}

The binding term carries the one parity-odd ingredient of the interior, the sign of $\kappa$, and it sits entirely in a total derivative.
Since $\bE\cdot\bB$ is a single-valued function on the configuration space, the canonical transformation $p\mapsto p+2m\kappa\,\eta GX$ removes it,
classically and as a phase of wave functions. The structure is that of the $\theta$-term of electrodynamics, a parity-odd coefficient times a
total derivative, with the difference that here it has no topological sector in which to act. In the positive form the same cross term is
gyroscopic, with Euler--Lagrange expression $2m\kappa G\dot X$, and the sign of $\kappa$ enters the equations of motion; that is another way of
seeing why the two forms behave so differently.

\subsection{What the binding root selects, and what tears it}

The binding root is not merely available; it carries structure. By \cref{thm:GLorentz} its centraliser is a Lorentz algebra $\lG$, not
conjugate to the interior one (\cref{prop:notconj}). The same operator limits what transport can do to the bound sector.

\begin{proposition}[Transport and a homogeneous threshold]
\label{prop:tear}
Let a connection along the path take values $\omega$ in $\so{3,3}$. If $\omega$ lies in the compact part of $\lG$, transport and the shift commute
and the orbit is carried rigidly. If $\omega$ has a boost component of strength $c$ in a pairing plane $(e_a,e_{a+3})$, the homogeneous $2\times2$
block there has exponents $\pm\sqrt{c^2-\kappa^2}$: oscillatory for $c<\kappa$ and real for $c>\kappa$.
\end{proposition}

\begin{proof}
In the basis $(e_a,e_{a+3})$ the boost generator is $\sigma=\bigl(\begin{smallmatrix}0&1\\1&0\end{smallmatrix}\bigr)$ and $G$ is
$\bigl(\begin{smallmatrix}0&-1\\1&0\end{smallmatrix}\bigr)$, so $-c\sigma+\kappa G=\bigl(\begin{smallmatrix}0&-c-\kappa\\-c+\kappa&0\end{smallmatrix}\bigr)$, whose
square is $(c^2-\kappa^2)\id$.
\end{proof}

The threshold classifies exponents of the homogeneous block, not boundedness of the full flow, and the difference is real. With the momentum
included, the phase-space generator $M=\bigl(\begin{smallmatrix}-c\sigma-\kappa G&\id/m\\0&-c\sigma+\kappa G\end{smallmatrix}\bigr)$ satisfies
$M^2-(c^2-\kappa^2)\id=\bigl(\begin{smallmatrix}0&-2c\sigma/m\\0&0\end{smallmatrix}\bigr)$, nilpotent and nonzero for $c\neq0$. At $\kappa=m=1$ and
$c=\tfrac12$ one finds $\det(\lambda\id-M)=(\lambda^2+\tfrac34)^2$ with $\rank(M^2+\tfrac34\id)=2$: Jordan blocks, and generic secular growth strictly
below the threshold. Along an open path, moreover, a connection coefficient can be changed by a frame rotation, so growth on one worldline
proves nothing; the invariant is the holonomy around a closed path. \Cref{sec:curved} computes it in an explicit curved model and finds that
normal curvature converts zero-momentum orbits into momentum-carrying ones, through exactly the failure of the normal rotation to commute with
$G$.

What this part delivers is a bound sector, a frequency, a centralising Lorentz algebra and an exponent threshold, all from an adopted binding
operator. It does not deliver a spectrum: $\kappa$ is not fixed, the single frequency is degenerate across the three pairs, and the energy is
indefinite, so these are bound orbits and not bound states in any quantum sense. Nor does it settle the relation between the two square roots of
the pairing beyond \cref{lem:Gclifford}: $\cK$ is $\eta$-antisymmetric and grades, $G$ is $\eta$-self-adjoint and binds, and they differ by the
Cartan involution $\Ad_{e_1e_2e_3}$.
